% SOURCE docs/release/2026-10-29/source/docs/writing-repair/2026-10-29/src/3f78a17ac47981c2.tex CHAR 0
\chapter{Design-based estimation, weak information and diagnostics}
\label{sw:design}

The preceding comparisons establish why a common interface must preserve
design-specific meanings. This chapter gathers their consequences for the
analyst's report: state the instrument claim, retain the complete uncertainty
set, distinguish threshold designs and show what a sensitivity parameter does
and does not reveal. The detailed implementation findings remain in the
preceding chapter; the purpose here is to connect them to the bank's decision.

\section{An instrument is a scientific claim}

Suppose a bank uses a rule-based funding incentive \(Z\) as an instrument for
an offered rate \(A\), with retained balances \(Y\) as the outcome and baseline
controls \(W\). Software can encode a regression system, but it cannot establish
that the incentive has no direct effect on balances or that its assignment is
independent of omitted balance determinants. Those are properties of the
economic setting and information process. The package comparison begins after
the instrument, target population, timing and exclusion claim have been stated.

The selected \texttt{ivreg} documentation describes ordinary two-stage least
squares and robust M/MM regression choices. Exogenous regressors in the outcome
equation must also appear among the instruments. An adapter should therefore
retain explicit endogenous, exogenous and instrument column roles rather than
constructing the formula from a single list of predictors. The robust regression
option is an estimator choice; it does not supply a weak-instrument confidence
guarantee. The inspected package is ivreg 0.6-8. The separate AER implementation
is compared in Section~\ref{exec22:riv-inference}; shared function names do not
make the engines interchangeable. That bounded review does not qualify the
whole AER family. \citep{v12_sw_ivreg}

Weak information changes the result contract. The inspected \texttt{ivmodel}
documentation exposes Anderson--Rubin and conditional-likelihood-ratio routines
whose confidence-set outputs are matrices of interval components. A client
that extracts one lower and one upper endpoint can lose disconnectedness or
unboundedness. The appropriate object is
\[
 C=\bigcup_{j=1}^{J} I_j,
\]
where each component \(I_j\) retains its endpoints, possible infinite bounds and
method-specific interpretation. The number of components is part of the result,
not formatting metadata. \citep{v12_sw_ivmodel}

For illustration, a hypothetical reported set
\(C=(-\infty,-2]\cup[3,\infty)\) excludes zero while remaining unbounded in both
directions. Its interval hull is the whole real line. Replacing the set with
that hull changes the question ``is zero in the set?'' and replacing it with a
midpoint invents a finite estimate. Conversely, an empty inversion result is
not equivalent to a missing calculation. The proposed result schema must
preserve these distinctions before any banking decision rule uses the output.
This hypothetical set is not a computed result from ivmodel.

The two R packages are therefore complementary candidates. ivreg supplies a
regression interface with specific estimation and diagnostic choices; ivmodel
supplies selected weak-information test/inversion interfaces. A matched trial
must align instruments, controls, scaling, sampling assumptions and covariance
choices. It must retain weak and zero first stages, singular designs and
unbounded sets. A regression returning quickly on a well-conditioned example
would not qualify the weak-information case, and the presence of an inversion
API would not establish coverage for dependent bank panels.
The source-level CLR failure substitution described in
Section~\ref{exec22:riv-inference} also remains a condition on reuse:
a numerical failure must not be reported as scientific rejection merely
because the returned p-value is zero.

\section{Regression discontinuity is a family of designs}

A credit policy can change at a score cutoff, but several causal questions can
be attached to that threshold. A sharp design studies a deterministic assignment
change. A fuzzy design uses a discontinuity in treatment probability. A kink
design studies a change in slope. These require different estimands and
denominators. Treating them as one generic local regression would make an
apparently comparable benchmark compare different scientific problems.

The inspected rdrobust 4.0.0 documentation exposes sharp, fuzzy and kink
settings, local-polynomial choices, bias correction, bandwidth choices and
variance options. The adapter should bind the cutoff and running-variable
units to the result and export the actual bandwidth and inference settings.
The bandwidth for a point-estimation objective and the choices underlying a
bias-corrected interval cannot be silently conflated. The saved 4.0.0 snapshot
predates the literature cutoff. Section~\ref{exec22:riv-scope} explains why
the subsequently authenticated CRAN publication time retains that distribution
for the client cutoff; it does not settle every upstream channel's history.
\citep{v12_sw_rdrobust}

For a fuzzy design, the treatment jump matters as much as the outcome jump.
If a local effect is expressed schematically as \(\tau_Y/\tau_A\), a small
\(\tau_A\) produces an unstable ratio even when each fitted regression has a
clean numerical convergence report. The bank's first question should be
whether the score rule changes the specified treatment in the relevant
population. The software trial should vary that first-stage jump and preserve
weak cases rather than selecting only a strong demonstration.

rddensity supplies a different operation: local-polynomial density or
manipulation diagnostics near the cutoff. Such diagnostics can challenge a
design, but failure to reject a discontinuity does not verify the assignment
mechanism or continuity of potential outcomes. An interface that displays
``RD valid'' after a density test would convert one diagnostic into a much
stronger scientific statement. A useful report instead gives the test,
bandwidth, support and interpretation limit alongside other design evidence.
Discrete scores and mass points require explicit attention. \citep{v12_sw_rddensity}

rdlocrand represents another approach, based on a declared local-randomization
window. Its selected documentation states that a positive polynomial order
uses asymptotic inference and that confidence-set inversion is performed over
a supplied grid. Thus a finite grid endpoint is not automatically a finite
confidence bound. If all evaluated values near the boundary remain accepted,
the result should say that the search is truncated. Extending the grid is a
computational remedy; it is not evidence for the local-randomization assumption.
Window selection, assignment probabilities and the intervention must still be
recorded. \citep{v12_sw_rdlocrand}

These packages should be evaluated as a collection of design-specific
operations. The proposed integration uses a common report for the score,
cutoff, assignment rule and target, then separate method records. The product
should not expose a single switch called ``robust RD'' that hides whether the
calculation is continuity-based local-polynomial inference or local-randomization
inference. The alternatives may both be useful, but they answer their questions
under different assumptions.

\section{Sensitivity parameters are not missing truth}

Two other R interfaces provide useful sensitivity summaries. sensemakr
parameterizes omitted-variable sensitivity through partial \(R^2\) quantities
and observed-covariate benchmarks. A large benchmark asks how conclusions would
change if an omitted variable had specified residual associations; it does not
measure the actual association of an unobserved variable. The bank must identify
what makes a proposed benchmark plausible, and the report should retain both
the chosen benchmark and the effect scale. \citep{v12_sw_sensemakr}

The selected EValue \texttt{evalues.RR} interface takes a risk ratio and optional
confidence limits and reports E-values for the point and the near-null confidence
limit. This is a particular scale-specific sensitivity operation. A default
probability difference, a standardized regression coefficient and a net present
value should not be passed to it as if they were risk ratios. Any conversion
needs an explicit target and assumptions. An E-value does not estimate a
confounding correction or certify an absence of confounding. \citep{v12_sw_evalue}

These tools can make a funding or credit recommendation more intelligible:
they expose how much an assumption would need to change before a conclusion
changes. They cannot supply verified instrument exclusion, known ownership,
correct record inclusion or a gold audit label. Their adoption is therefore
conditional on a result contract that separates observed data summaries,
user-supplied sensitivity assumptions and causal conclusions.


